% Chapter 1

\chapter{Introduction} % Main chapter title

\label{Chapter1} % For referencing the chapter elsewhere, use \ref{Chapter1}

%----------------------------------------------------------------------------------------

Each academic semester, the university faces the same challenging problem. How do you allocate hundreds of classes, a number of lecturers, and a fixed set of class rooms to a semester-wide timetable? This problem is known as the University Timetabling Problem (UTP) and it has been proven NP-hard, i.e., that there is no efficient algorithm for solving the problem optimally. This chapter will give a brief overview of the IIITA Timetabling System, describe its rationale, and outline the remaining chapters in this thesis.
%-----------------------------------
\section{Background and Motivation}

An academic timetable involves assigning lectures, tutorials, and laboratory classes to a particular time slot, venue, and lecturer, while taking into account the limitations imposed by the situation. There are some constraints that are rigid, in that it is not possible to double book the lecturer or run two classes in one room at the same time. On the other hand, there are certain soft constraints which include having the minimum amount of gap between lectures, and the professors prefer lecturing in one half of the day only.

IIIT Allahabad is one among the many colleges in India where the scheduling process remains manual. Timetables are created manually using Microsoft Excel, and conflicts are manually resolved by the academic coordinators. Even though Microsoft Excel is highly customizable and flexible, it does not have any sense of the meaning of a ``professor'' or ``room''. It is up to the manual efforts of the people preparing the timetables to make sure that such problems are avoided.

The issues do not stop even after the timetable is released. Using a static PDF or Excel file will result in the need to send out the whole document every time there is a change, such as a classroom change or an instructor change. The students will have to work with outdated versions, and no one will be able to answer basic queries such as "What classrooms are available at 2 PM on Wednesday?"

These annoyances were what pushed us to develop the IIITA Timetabling System. We sought a unified solution that would become the ultimate authority on all scheduling information. This system should be capable of automatically generating conflict-free timetables, allowing the administrators to adjust them interactively, and providing individualized timetables for each stakeholder.


%-----------------------------------
\section{Problem Statement}

The underlying issue we sought to address is simple: The manual creation and management of timetables via spreadsheets at IIITA is inefficient, error-prone, and inherently inflexible. Specifically, the existing process suffers from:

\begin{itemize}
    \item \textbf{Unstructured Data:} Our schedule is created in excel sheets with merged cells, non-matching names, and relations that are beyond any computer program’s comprehension.
    \item \textbf{Manual Conflict Resolution:} To check whether there are problems like clash of classrooms, unavailability of certain teachers or courses, one needs to do a whole process manually, hence increasing the possibility of error.
    \item \textbf{No Automated Generation:}Currently, there is no system to automatically create a non-conflict timetable according to the needs of our institution.
    \item \textbf{One-Size-Fits-All Views:} It makes no sense to make a timetable which will not satisfy everyone. Students need to search in this timetable, and teachers cannot quickly find their schedules.
    \item \textbf{Resource Management:} Administrators lack a simple way of knowing which classrooms are vacant at certain times and how effectively our resources are being used.
    \item \textbf{Static Distribution:} After publishing the timetable in pdf form, updating it means distributing it all over again.
\end{itemize}

Thus, our objective was to develop a system which could help close the chasm between unstructured spreadsheet formats and structured data, schedule creation, editing tools, and provide role-based visualization of the timetable.

%-----------------------------------
\section{Objectives of the Study}

We set ourselves seven concrete objectives:

\begin{enumerate}
    \item \textbf{Heuristic Data Ingestion:} Design an engine that will parse IIITA’s complicated Excel timetable, extract metadata (course subjects, instructors, class venues), merge cells, and normalise into a relational schema.
    \item \textbf{Real-Time Clash Detection:} Design a set of algorithms for detecting conflicts related to time, room availability, instructor workload, and student sections immediately upon importing timetable data.
    \item \textbf{Role-Based Views:} Design role-based UI interfaces for different actors in the process including:
    \begin{itemize}
        \item Student Interface with schedule filtering options and electives.
        \item Professor Interface with personal schedule only visible to the professor.
        \item Admin Interface with master schedule, free room finder, and other administrative tools.
    \end{itemize}
    \item \textbf{Export and Reports:} Provide users with options to export schedule as PDF, as well as providing Google Sheets support for administrative purposes.
    \item \textbf{Automated Timetable Generation:} Implement a $(1+1)$ Evolution Strategy solver that evaluates candidates against 10 hard constraints and 3 soft constraints, adapts its step size via the 1/5th success rule, supports warm-starting from existing schedules, and handling inter-semester constraints through locked sessions.
    \item \textbf{Interactive Timetable Editor:} Develop a graphical drag-and-drop editor for manual timetable optimisation with instant conflict detection, room and professor modification capabilities, and one-click Large Neightbourhood Search (LNS) auto-resolution for any remaining conflicts.
    \item \textbf{Intelligent Generator:} Create a guided, multi-step workflow that guides administrators through cluster selection, home room mapping, professor assignment, and solver configuration.
\end{enumerate}

%-----------------------------------
\section{Scope and Limitations}

The scope of the project is to implement all phases of the IIITA Timetabling System, including the React + TypeScript front end, the PostgreSQL back end with Supabase, and the browser-based execution of the (1+1)-ES. Even though the system has been designed specifically for the timetable of IIIT Allahabad, it is still extensible for other colleges.

That said, there are a few important limitations to keep in mind:
\begin{itemize}
    \item \textbf{Parser Sensitivity:} Our Excel parser relies on heuristic rules tuned to IIITA's template. Major deviations from that format may require tweaking the parsing logic.
    \item \textbf{Single-Threaded Solver:} The solver runs on the browser's main thread using batched \texttt{setTimeout} calls. This strategy will work well with the present input size, though in case of very large inputs it may become a performance bottleneck.
    \item \textbf{Heuristic Optimality:} Evolution Strategies are meta-heuristics. They guarantee feasibility by providing solutions with zero hard constraint violations, though there is no assurance of optimality for soft constraints based solutions. This depends upon the number of generations.
\end{itemize}

%-----------------------------------
\section{Organization of the Thesis}

The rest of this thesis is organised as follows:

\textbf{Chapter 2: Literature Review} will provide a review of current researches into university timetabling, including the basic theories of UTP as an NP-hard problem and current algorithmic solutions in literature.

\textbf{Chapter 3: Technologies and Frameworks} will present our choice of technology stack for this project, i.e., React 19, TypeScript, Tailwind CSS, Vite, and Supabase, along with the reasons for their adoption.

\textbf{Chapter 4: System Architecture and Design} will detail the database schema, the data ingest pipeline, the solver architecture, and finally, the client-side routing of components.

\textbf{Chapter 5: Implementation and Results} will form the core of this thesis, where the implementation details of the Excel parser, clash detection algorithms, (1+1)-ES solver and constraint system, drag-and-drop editor and auto-LNS fix, and the intelligent generator will be described.

\textbf{Chapter 6: Conclusion and Future Work} sums up what we achieved, reflects on the challenges we faced, and outlines future directions such as parallelised solvers and hybrid meta-heuristics.
